% GENERATED FILE. Edit canonical lessons, then run npm run generate:books.
% canonical-source-hash: fnv1a64:2552b7e403a181d3
% canonical-lessons: RU-C524-detektiv, RU-C524-fantastika, RU-C524-biografiya, RU-C524-povest, RU-C524-memuary, RU-R524-review-517-524

\chapter{Kinds of Books}
\label{ch:ru-kinds-of-books}
\hlchaptermodality{524}

\begin{chapteropening}
\textbf{By the end of this chapter:} \emph{I can say what kinds of books I like.}

Complete the last lesson of chapter 524: I can say what kinds of books I like.
\end{chapteropening}

\section[\texorpdfstring{\ru{детектив} (detektív)}{детектив (detektív)}]{\ru{детектив} (detektív) — a detective story; a detective}

\label{lesson:RU-C524-detektiv}



\begin{quote}
Before the new one: say the Russian for an ending, then the Russian for to get engrossed in reading.
\end{quote}

\subsection*{You'll want to know: \ru{детектив}}
\textbf{\ru{детектив}} (\emph{detektív}) — \textquotedblleft{}a detective story; a detective\textquotedblright{}.

\begin{quote}
\ru{Летом} \ru{я} \ru{люблю} \ru{читать} \ru{детективы}.
\end{quote}

\begin{quote}
\emph{Létom ya lyublyú chitát detektívy.}
\end{quote}

\begin{quote}
In the summer I like reading detective stories.
\end{quote}

Masculine.

\begin{cousinweb}
EtymDB lists \emph{\ru{детектив}} without a source; it is usually traced to English \emph{detective}.
\end{cousinweb}

\begin{grammarlens}[title={What you've built}]
One more word for telling, explaining and describing at B1.
\end{grammarlens}

\subsection*{Your turn}
\begin{itemize}
\raggedright
  \item \emph{Say it:} \emph{detektív}
  \item \emph{Say it:} \emph{detektív}, once more
  \item \emph{Say it:} a sentence of your own with \emph{detektív}
  \item \emph{Recall:} say the Russian for an ending, then the Russian for to get engrossed in reading, then say \emph{detektív} again
\end{itemize}

\begin{tcolorbox}[breakable,colback=teal!4,colframe=teal!35!black,title={Before you move on}]
What does \ru{детектив} mean? (\textquotedblleft{}A detective story; a detective\textquotedblright{}.) Say it once more.
\end{tcolorbox}
\section[\texorpdfstring{\ru{фантастика} (fantástika)}{фантастика (fantástika)}]{\ru{фантастика} (fantástika) — science fiction}

\label{lesson:RU-C524-fantastika}



\begin{quote}
Before the new one: say the Russian for to get engrossed in reading, then the Russian for a detective story.
\end{quote}

\subsection*{You'll want to know: \ru{фантастика}}
\textbf{\ru{фантастика}} (\emph{fantástika}) — \textquotedblleft{}science fiction\textquotedblright{}.

\begin{quote}
\ru{Мой} \ru{брат} \ru{читает} \ru{только} \ru{фантастику}.
\end{quote}

\begin{quote}
\emph{Moy brat chitáyet tólko fantástiku.}
\end{quote}

\begin{quote}
My brother only reads science fiction.
\end{quote}

Feminine. Also fantasy, and colloquially something amazing: \emph{\ru{Это} \ru{фантастика}!}

\begin{cousinweb}
EtymDB lists \emph{\ru{фантастика}} without a source; it is usually traced through \emph{\ru{фантастический}} to Greek \emph{phantastikós} 'able to imagine'.
\end{cousinweb}

\begin{grammarlens}[title={What you've built}]
One more word for telling, explaining and describing at B1.
\end{grammarlens}

\subsection*{Your turn}
\begin{itemize}
\raggedright
  \item \emph{Say it:} \emph{fantástika}
  \item \emph{Say it:} \emph{fantástika}, once more
  \item \emph{Say it:} a sentence of your own with \emph{fantástika}
  \item \emph{Recall:} say the Russian for to get engrossed in reading, then the Russian for a detective story, then say \emph{fantástika} again
\end{itemize}

\begin{tcolorbox}[breakable,colback=teal!4,colframe=teal!35!black,title={Before you move on}]
What does \ru{фантастика} mean? (\textquotedblleft{}Science fiction\textquotedblright{}.) Say it once more.
\end{tcolorbox}
\section[\texorpdfstring{\ru{биография} (biográfiya)}{биография (biográfiya)}]{\ru{биография} (biográfiya) — a biography}

\label{lesson:RU-C524-biografiya}



\begin{quote}
Before the new one: say the Russian for a detective story, then the Russian for science fiction.
\end{quote}

\subsection*{You'll want to know: \ru{биография}}
\textbf{\ru{биография}} (\emph{biográfiya}) — \textquotedblleft{}a biography\textquotedblright{}.

\begin{quote}
\ru{Я} \ru{прочитал} \ru{биографию} \ru{известного} \ru{учёного}.
\end{quote}

\begin{quote}
\emph{Ya prochitál biográfiyu izvéstnogo uchyónogo.}
\end{quote}

\begin{quote}
I've read the biography of a famous scientist.
\end{quote}

Feminine.

\begin{cousinweb}
\emph{\ru{биография}} is a learned word built from Greek \emph{bíos} 'life' and \emph{graph\'{\={e}}} 'writing'; EtymDB lists the Russian word without a source.
\end{cousinweb}

\begin{grammarlens}[title={What you've built}]
One more word for telling, explaining and describing at B1.
\end{grammarlens}

\subsection*{Your turn}
\begin{itemize}
\raggedright
  \item \emph{Say it:} \emph{biográfiya}
  \item \emph{Say it:} \emph{biográfiya}, once more
  \item \emph{Say it:} a sentence of your own with \emph{biográfiya}
  \item \emph{Recall:} say the Russian for a detective story, then the Russian for science fiction, then say \emph{biográfiya} again
\end{itemize}

\begin{tcolorbox}[breakable,colback=teal!4,colframe=teal!35!black,title={Before you move on}]
What does \ru{биография} mean? (\textquotedblleft{}A biography\textquotedblright{}.) Say it once more.
\end{tcolorbox}
\section[\texorpdfstring{\ru{повесть} (póvest)}{повесть (póvest)}]{\ru{повесть} (póvest) — a novella}

\label{lesson:RU-C524-povest}



\begin{quote}
Before the new one: say the Russian for science fiction, then the Russian for a biography.
\end{quote}

\subsection*{You'll want to know: \ru{повесть}}
\textbf{\ru{повесть}} (\emph{póvest}) — \textquotedblleft{}a novella\textquotedblright{}.

\begin{quote}
\ru{Эту} \ru{повесть} \ru{мы} \ru{читали} \ru{в} \ru{школе}.
\end{quote}

\begin{quote}
\emph{Étu póvest my chitáli v shkóle.}
\end{quote}

\begin{quote}
We read this novella at school.
\end{quote}

Feminine. Longer than \emph{\ru{рассказ}}, shorter than \emph{\ru{роман}}.

\begin{cousinweb}
\emph{\ru{повесть}} is built on \emph{\ru{по}-} and \emph{\ru{весть}} 'news, tidings', so a story is something told; EtymDB lists \emph{\ru{повесть}} without a source.
\end{cousinweb}

\begin{grammarlens}[title={What you've built}]
One more word for telling, explaining and describing at B1.
\end{grammarlens}

\subsection*{Your turn}
\begin{itemize}
\raggedright
  \item \emph{Say it:} \emph{póvest}
  \item \emph{Say it:} \emph{póvest}, once more
  \item \emph{Say it:} a sentence of your own with \emph{póvest}
  \item \emph{Recall:} say the Russian for science fiction, then the Russian for a biography, then say \emph{póvest} again
\end{itemize}

\begin{tcolorbox}[breakable,colback=teal!4,colframe=teal!35!black,title={Before you move on}]
What does \ru{повесть} mean? (\textquotedblleft{}A novella\textquotedblright{}.) Say it once more.
\end{tcolorbox}
\section[\texorpdfstring{\ru{мемуары} (memuáry)}{мемуары (memuáry)}]{\ru{мемуары} (memuáry) — memoirs}

\label{lesson:RU-C524-memuary}



\begin{quote}
Before the new one: say the Russian for a biography, then the Russian for a novella.
\end{quote}

\subsection*{You'll want to know: \ru{мемуары}}
\textbf{\ru{мемуары}} (\emph{memuáry}) — \textquotedblleft{}memoirs\textquotedblright{}.

\begin{quote}
\ru{Старый} \ru{актёр} \ru{написал} \ru{мемуары}.
\end{quote}

\begin{quote}
\emph{Stáryy aktyór napisál memuáry.}
\end{quote}

\begin{quote}
The old actor wrote his memoirs.
\end{quote}

Plural only.

\begin{cousinweb}
\emph{\ru{мемуары}} was borrowed from French \emph{mémoires} 'memoirs' (EtymDB), which is usually traced to Latin \emph{memoria} 'memory'.
\end{cousinweb}

\begin{grammarlens}[title={What you've built}]
One more word for telling, explaining and describing at B1.
\end{grammarlens}

\subsection*{Your turn}
\begin{itemize}
\raggedright
  \item \emph{Say it:} \emph{memuáry}
  \item \emph{Say it:} \emph{memuáry}, once more
  \item \emph{Say it:} a sentence of your own with \emph{memuáry}
  \item \emph{Recall:} say the Russian for a biography, then the Russian for a novella, then say \emph{memuáry} again
\end{itemize}

\begin{tcolorbox}[breakable,colback=teal!4,colframe=teal!35!black,title={Before you move on}]
What does \ru{мемуары} mean? (\textquotedblleft{}Memoirs\textquotedblright{}.) Say it once more.
\end{tcolorbox}
\section[(dialogue)]{Words from chapters 517–524 — a review}

\label{lesson:RU-R524-review-517-524}



\begin{quote}
Say the Russian for a choir and an orchestra.
\end{quote}

\subsection*{Your turn}
Say these aloud:

\begin{itemize}
\raggedright
  \item a choir, an orchestra, a conductor, to conduct, a composer
  \item to perform, a performer, a rehearsal, to rehearse, a tour
  \item a stage curtain, an interval, to applaud, a premiere, to bow
  \item a director, an actress, a role, a viewer, to act in a film
  \item a paintbrush, canvas, watercolour, a landscape, to pose
  \item a sculptor, a sculpture, marble, a statue, to inspire
  \item a plot, a character, a heroine, an ending, to get engrossed in reading
  \item a detective story, science fiction, a biography, a novella, memoirs
\end{itemize}

\begin{tcolorbox}[breakable,colback=teal!4,colframe=teal!35!black,title={Before you move on}]
A choir? (\textbf{\ru{хор}}.) An orchestra? (\textbf{\ru{оркестр}}.) A conductor? (\textbf{\ru{дирижёр}}.) To conduct? (\textbf{\ru{дирижировать}}.) A composer? (\textbf{\ru{композитор}}.) To perform? (\textbf{\ru{исполнять}}.) A performer? (\textbf{\ru{исполнитель}}.) A rehearsal? (\textbf{\ru{репетиция}}.) To rehearse? (\textbf{\ru{репетировать}}.) A tour? (\textbf{\ru{гастроли}}.) A stage curtain? (\textbf{\ru{занавес}}.) An interval? (\textbf{\ru{антракт}}.) To applaud? (\textbf{\ru{аплодировать}}.) A premiere? (\textbf{\ru{премьера}}.) To bow? (\textbf{\ru{кланяться}}.) A director? (\textbf{\ru{режиссёр}}.) An actress? (\textbf{\ru{актриса}}.) A role? (\textbf{\ru{роль}}.) A viewer? (\textbf{\ru{зритель}}.) To act in a film? (\textbf{\ru{сниматься}}.) A paintbrush? (\textbf{\ru{кисть}}.) Canvas? (\textbf{\ru{холст}}.) Watercolour? (\textbf{\ru{акварель}}.) A landscape? (\textbf{\ru{пейзаж}}.) To pose? (\textbf{\ru{позировать}}.) A sculptor? (\textbf{\ru{скульптор}}.) A sculpture? (\textbf{\ru{скульптура}}.) Marble? (\textbf{\ru{мрамор}}.) A statue? (\textbf{\ru{статуя}}.) To inspire? (\textbf{\ru{вдохновлять}}.) A plot? (\textbf{\ru{сюжет}}.) A character? (\textbf{\ru{персонаж}}.) A heroine? (\textbf{\ru{героиня}}.) An ending? (\textbf{\ru{развязка}}.) To get engrossed in reading? (\textbf{\ru{зачитываться}}.) A detective story? (\textbf{\ru{детектив}}.) Science fiction? (\textbf{\ru{фантастика}}.) A biography? (\textbf{\ru{биография}}.) A novella? (\textbf{\ru{повесть}}.) Memoirs? (\textbf{\ru{мемуары}}.)
\end{tcolorbox}
